\chapter{Stroomkringen: lussen, serie, parallel}\label{ch:g7:circuits-series-parallel}

Vier jaar schakelingen hebben je gereedschapskist gevuld: lussen,
schema's, de twee grote bedradingen. Maar één personage heeft al die
tijd naamloos gewerkt --- het iets dat stroomt. Dit jaar krijgt het
zijn naam, zijn richting en zijn vreemdste eigenschap: het raakt nooit
op. Met de stroom op het toneel houden serie en \omterm{def:g5:series-parallel-circuits:parallel}{parallel} op recepten
te zijn en worden ze wetten.

\section{De stroom betreedt het toneel}

\begin{definition}[Elektrische stroom]\label{def:g7:circuits-series-parallel:current}
De \emph{elektrische stroom}\index{elektrische stroom} is de stroming
die rondgaat in een \omterm{def:g3:first-electric-circuit:circuit}{gesloten stroomkring}: een menigte geladen
deeltjes --- onvoorstelbaar talrijk, onvoorstelbaar klein, neven van
de moleculaire wereld --- die door de duw van de \omterm{def:g3:first-electric-circuit:battery}{batterij} door het
metaal van de draden aan het marcheren worden gezet. Geen lus, geen
mars: de stroom bestaat alleen zolang de weg gesloten is. (Wat de
marcheerders precies zijn, en hoe een metaal ze doorlaat, is een mooi
verhaal dat voor de laatste jaren van dit boek bewaard blijft.)
\end{definition}

\begin{proposition}[De afgesproken richting]\label{prop:g7:circuits-series-parallel:direction}
Volgens wereldwijde afspraak tekenen \omterm{def:g6:circuits-and-safety:diagram}{schakelschema}'s de stroom als
lopend \emph{buiten de \omterm{def:g3:first-electric-circuit:battery}{batterij} van de $+$pool, rond de kring,
naar de $-$pool}. Elke pijl op elk schema ter wereld gehoorzaamt
aan die afspraak --- die gemaakt werd lang voordat iemand het de
marcheerders zelf kon vragen. (Toen de natuurkunde het eindelijk kon
vragen, zat er een beroemd grapje in het antwoord, bewaard voor een
later jaar.)
\end{proposition}

\begin{omfigure}
\begin{tikzpicture}[scale=0.9]
  \draw (0,0) to[battery1, l=batterij] (0,3.0)
    to[short] (2.0,3.0)
    to[lamp] (4.4,3.0)
    to[short] (6.4,3.0)
    to[short] (6.4,0)
    to[short] (0,0);
  \draw[-Stealth, very thick, omThm] (1.1,3.25) -- (1.9,3.25);
  \draw[-Stealth, very thick, omThm] (6.65,1.9) -- (6.65,1.1);
  \draw[-Stealth, very thick, omThm] (5.3,-0.25) -- (4.5,-0.25);
  \draw[-Stealth, very thick, omThm] (1.9,-0.25) -- (1.1,-0.25);
  \node[right, font=\small, omThm] at (7.0,1.5)
    {de stroom: uit $+$, rond, naar $-$};
\end{tikzpicture}

{\small De afspraak die elke elektricien op Aarde deelt: de
stroompijlen verlaten de $+$ van de \omterm{def:g3:first-electric-circuit:battery}{batterij}, doen de lus rond en
keren terug via $-$.}
\end{omfigure}

\begin{proposition}[De stroom wordt niet verbruikt]\label{prop:g7:circuits-series-parallel:notconsumed}
Apparaten eten de stroom niet op. Alle stroom die een lamp,
\omterm{ex:g6:circuits-and-safety:motor}{motor} of \omterm{ex:g6:circuits-and-safety:motor}{zoemer} binnengaat, verlaat hem volledig en marcheert
verder: wat het apparaat neemt is \emph{\omterm{def:g5:everyday-energy:energy}{energie}} uit de voorraad van
de \omterm{def:g3:first-electric-circuit:battery}{batterij}, bezorgd door de langstrekkende mars. In elke serielus
loopt daarom één en dezelfde stroom door elk apparaat --- er komt
evenveel terug bij de \omterm{def:g3:first-electric-circuit:battery}{batterij} als er ooit vertrok.
\end{proposition}

\begin{example}[De misvatting voor het gerecht]\label{ex:g7:circuits-series-parallel:trial}
De verleidelijke fout: ``het eerste \omterm{def:g3:first-electric-circuit:bulb}{lampje} verbruikt wat stroom, dus
moet het tweede zwakker gloeien.'' De kring weerlegt haar zelf. Zet
twee \emph{identieke} \omterm{def:g3:first-electric-circuit:bulb}{lampjes} \omterm{def:g4:series-circuits:series}{in serie} en kijk: allebei gloeien ze
precies even fel --- en ze verwisselen verandert niets. Als er onderweg
stroom verbruikt werd, zou \omterm{def:g3:first-electric-circuit:bulb}{lampje} twee altijd het zwakke moeten zijn,
welk \omterm{def:g3:first-electric-circuit:bulb}{lampje} daar ook zit. Dat is het nooit. De stroom komt heel aan;
wat het paar deelt is de \emph{duw} van de \omterm{def:g3:first-electric-circuit:battery}{batterij}, en daarom zijn ze
allebei zwakker dan één enkel \omterm{def:g3:first-electric-circuit:bulb}{lampje} --- even zwak.
\end{example}

\section{De twee families, als wetten}

\begin{proposition}[De seriewet]\label{prop:g7:circuits-series-parallel:series}
In een \omterm{def:g4:series-circuits:series}{serieschakeling} --- één enkele lus --- loopt door elk apparaat
dezelfde stroom: één weg, één mars, geen uitzonderingen. Gevolgen, nu
verklaard: één onderbreking legt de mars overal tegelijk stil;
identieke apparaten delen het werk gelijkelijk, waar ze ook zitten; en
de volgorde van de apparaten op de lus kan niets uitmaken voor een
mars die door alle apparaten even goed gaat.
\end{proposition}

\begin{proposition}[De parallelwet]\label{prop:g7:circuits-series-parallel:parallel}
Bij een \omterm{def:g5:series-parallel-circuits:parallel}{knooppunt} \emph{splitst} de stroom zich: een deel van de mars
neemt elke \omterm{def:g5:series-parallel-circuits:parallel}{tak}, en de delen voegen zich compleet weer samen bij het
verre \omterm{def:g5:series-parallel-circuits:parallel}{knooppunt} --- niets verloren, niets gewonnen. Elke \omterm{def:g5:series-parallel-circuits:parallel}{tak}
ontvangt de volle duw van de \omterm{def:g3:first-electric-circuit:battery}{batterij}, dus werkt elk apparaat alsof
het alleen was; en een geblokkeerde \omterm{def:g5:series-parallel-circuits:parallel}{tak} stuurt eenvoudig niemand,
terwijl de andere \omterm{def:g5:series-parallel-circuits:parallel}{takken} doormarcheren.
\end{proposition}

\begin{omfigure}
\begin{tikzpicture}[scale=0.85]
  \draw (0,0) to[battery1] (0,3.2)
    to[short] (2.4,3.2);
  \draw (2.4,3.2) to[short] (2.4,2.3)
    to[lamp] (5.6,2.3) to[short] (5.6,3.2);
  \draw (2.4,3.2) to[short] (2.4,4.1)
    to[lamp] (5.6,4.1) to[short] (5.6,3.2);
  \draw (5.6,3.2) to[short] (7.2,3.2) to[short] (7.2,0)
    to[short] (0,0);
  \fill (2.4,3.2) circle (2pt);
  \fill (5.6,3.2) circle (2pt);
  \draw[-Stealth, very thick, omThm] (1.0,3.45) -- (1.9,3.45);
  \draw[-Stealth, thick, omThm] (2.85,4.35) -- (3.65,4.35);
  \draw[-Stealth, thick, omThm] (2.85,2.05) -- (3.65,2.05);
  \draw[-Stealth, very thick, omThm] (6.05,3.45) -- (6.95,3.45);
  \node[above, font=\small, omThm] at (1.45,3.55) {hele mars};
  \node[above, font=\small, omThm] at (6.5,3.55) {weer heel};
\end{tikzpicture}

{\small De parallelwet getekend: de mars splitst bij het eerste
\omterm{def:g5:series-parallel-circuits:parallel}{knooppunt}, elk deel doet zijn \omterm{def:g5:series-parallel-circuits:parallel}{tak}, en de volle mars stelt zich bij
het tweede weer samen.}
\end{omfigure}

\begin{example}[Felheid lezen als een elektricien]\label{ex:g7:circuits-series-parallel:brightness}
Eén \omterm{def:g3:first-electric-circuit:battery}{batterij}, identieke \omterm{def:g3:first-electric-circuit:bulb}{lampjes}. Alleen: volle felheid --- de
referentie. Twee \omterm{def:g4:series-circuits:series}{in serie}: allebei even zwak --- één mars, gedeelde
duw. Twee \omterm{def:g5:series-parallel-circuits:parallel}{parallel}: allebei volop fel --- elke \omterm{def:g5:series-parallel-circuits:parallel}{tak} geniet van de
hele duw (en de \omterm{def:g3:first-electric-circuit:battery}{batterij} raakt ruwweg twee keer zo snel leeg: niets
is gratis). Drie \omterm{def:g4:series-circuits:series}{in serie}: nog zwakker, allemaal gelijk. De blik van
de elektricien: tel de lusgenoten voor de zwakte, tel de
\omterm{def:g5:series-parallel-circuits:parallel}{takken} voor de batterijrekening.
\end{example}

\begin{method}[Het lot van elke lamp voorspellen]\label{met:g7:circuits-series-parallel:predict}
Bij een gegeven schema, voor elke lamp:
\begin{enumerate}
  \item \emph{Bestaan}: is er minstens één ononderbroken lus van $+$
        door deze lamp naar $-$? Geen lus, geen \omterm{def:g1:light-and-shadows:source}{licht}.
  \item \emph{Gezelschap op de weg}: welke apparaten delen
        \emph{elke} lus door deze lamp? Dat zijn haar seriegenoten
        --- meer genoten, zwakkere gloed.
  \item \emph{Buren op andere \omterm{def:g5:series-parallel-circuits:parallel}{takken}}: welke apparaten zitten op
        andere \omterm{def:g5:series-parallel-circuits:parallel}{takken}? Die maken deze lamp niet zwakker en gaan er
        ook niet mee uit.
  \item \emph{\omterm{ex:g3:first-electric-circuit:switch}{Schakelaars}}: een open \omterm{ex:g3:first-electric-circuit:switch}{schakelaar} doodt precies de
        lussen waarop hij zit --- doe stap 1 opnieuw met die lussen
        weggedacht.
\end{enumerate}
\end{method}

\begin{remark}[Waarom de misvatting zo hardnekkig is]\label{rem:g7:circuits-series-parallel:whyhard}
``Er moet toch iets opgaan --- de \omterm{def:g3:first-electric-circuit:battery}{batterij} raakt leeg!'' Waar: de
opgeslagen \emph{\omterm{def:g5:everyday-energy:energy}{energie}} van de \omterm{def:g3:first-electric-circuit:battery}{batterij} wordt uitgegeven en door
de langstrekkende mars als \omterm{def:g1:light-and-shadows:source}{licht} en \omterm{def:g5:heat-and-insulation:heat}{warmte} aan lampen overhandigd.
Maar de marcheerders zelf circuleren ongedeerd, als de stoeltjes van
een skilift die de hele dag skiërs de berg op brengen zonder dat er
ooit een stoeltje wordt opgegeten. Houd de twee boekhoudingen uit
elkaar --- stroom circuleert, \omterm{def:g5:everyday-energy:energy}{energie} wordt bezorgd --- en elke
schakeling in dit boek blijft eerlijk. Het instrument dat de mars zelf
telt, voegt zich volgend jaar bij ons.
\end{remark}

\section{Oefeningen}

\begin{exercise}[$\star$]\label{exo:g7:circuits-series-parallel:1}
Wat is de \omterm{def:g7:circuits-series-parallel:current}{elektrische stroom}? Wanneer bestaat hij in een
schakeling, en wanneer niet?
\end{exercise}

\begin{exercise}[$\star$]\label{exo:g7:circuits-series-parallel:2}
Noem de afgesproken richting van de stroom. Waar wijzen de pijlen op
een schema van een lus met één \omterm{def:g3:first-electric-circuit:bulb}{lampje}?
\end{exercise}

\begin{exercise}[$\star$]\label{exo:g7:circuits-series-parallel:3}
Twee identieke \omterm{def:g3:first-electric-circuit:bulb}{lampjes} \omterm{def:g4:series-circuits:series}{in serie} gloeien precies even fel. Welke
populaire fout weerlegt die gelijkheid, en hoe?
\end{exercise}

\begin{exercise}[$\star$]\label{exo:g7:circuits-series-parallel:4}
In een serielus met een \omterm{ex:g6:circuits-and-safety:motor}{motor} en een lamp wordt de stroom door de
\omterm{ex:g6:circuits-and-safety:motor}{motor} vergeleken met de stroom door de lamp. Wat zegt de seriewet?
\end{exercise}

\begin{exercise}[$\star$]\label{exo:g7:circuits-series-parallel:5}
Wat gebeurt er met de stroom bij het eerste \omterm{def:g5:series-parallel-circuits:parallel}{knooppunt} van een
\omterm{def:g5:series-parallel-circuits:parallel}{parallelschakeling}? En bij het tweede? Wat geldt er voor het
aandeel van elke \omterm{def:g5:series-parallel-circuits:parallel}{tak} in de duw van de \omterm{def:g3:first-electric-circuit:battery}{batterij}?
\end{exercise}

\begin{exercise}[$\star$]\label{exo:g7:circuits-series-parallel:6}
Eén \omterm{def:g3:first-electric-circuit:battery}{batterij}, identieke \omterm{def:g3:first-electric-circuit:bulb}{lampjes}: zet de felheid van een
\omterm{def:g3:first-electric-circuit:bulb}{lampje} op volgorde --- alleen; \ met één seriegenoot; \ met één
parallelbuur; \ met twee seriegenoten.
\end{exercise}

\begin{exercise}[$\star$]\label{exo:g7:circuits-series-parallel:7}
Welke opstelling maakt de \omterm{def:g3:first-electric-circuit:battery}{batterij} sneller leeg: twee \omterm{def:g3:first-electric-circuit:bulb}{lampjes}
\omterm{def:g4:series-circuits:series}{in serie}, of dezelfde twee \omterm{def:g5:series-parallel-circuits:parallel}{parallel}? Waarom, in marstaal?
\end{exercise}

\begin{exercise}[$\star$]\label{exo:g7:circuits-series-parallel:8}
Leg het skiliftbeeld van
\cref{rem:g7:circuits-series-parallel:whyhard} uit: wat speelt de
stoeltjes, wat speelt de skiërs, en wat wordt er werkelijk verbruikt?
\end{exercise}

\begin{exercise}[$\star\star$]\label{exo:g7:circuits-series-parallel:9}
Een schema: \omterm{def:g3:first-electric-circuit:battery}{batterij}, \omterm{ex:g3:first-electric-circuit:switch}{schakelaar}, dan lamp A; na A stuurt een
\omterm{def:g5:series-parallel-circuits:parallel}{knooppunt} twee \omterm{def:g5:series-parallel-circuits:parallel}{takken} weg --- lamp B op de ene, lamp C op de
andere --- die weer samenkomen vóór de terugkeer naar de
\omterm{def:g3:first-electric-circuit:battery}{batterij}. Met \cref{met:g7:circuits-series-parallel:predict}: wie zijn
de seriegenoten van A? Wat doet de gloed van A als de \omterm{def:g5:series-parallel-circuits:parallel}{tak} van C
wordt doorgeknipt?
\end{exercise}

\begin{exercise}[$\star\star$]\label{exo:g7:circuits-series-parallel:10}
Feestverlichting vraagt om een ontwerp: tien \omterm{def:g3:first-electric-circuit:bulb}{lampjes} die
allemaal op volle, onafhankelijke felheid schijnen uit één voeding ---
maar de hele set moet aan één hoofdschakelaar gehoorzamen. Teken of
beschrijf de bedrading, met een beroep op beide wetten.
\end{exercise}

\begin{exercise}[$\star\star$]\label{exo:g7:circuits-series-parallel:11}
In de \omterm{def:g5:series-parallel-circuits:parallel}{parallelschakeling} met twee \omterm{def:g3:first-electric-circuit:bulb}{lampjes} is de \omterm{def:g5:series-parallel-circuits:parallel}{tak} van
\omterm{def:g3:first-electric-circuit:bulb}{lampje} 1 dikke koperdraad, terwijl de \omterm{def:g5:series-parallel-circuits:parallel}{tak} van \omterm{def:g3:first-electric-circuit:bulb}{lampje} 2 een
lange, dunne draad bevat. \omterm{def:g3:first-electric-circuit:bulb}{Lampje} 2 gloeit een beetje zwakker. Wat
suggereert dat over hoe de weg van een \omterm{def:g5:series-parallel-circuits:parallel}{tak} zelf zijn aandeel in de
mars kan afknijpen? (Volgend jaar geeft dat afknijpen een naam en een
wet.)
\end{exercise}

\begin{exercise}[$\star\star\star$]\label{exo:g7:circuits-series-parallel:12}
Een \omterm{def:g3:first-electric-circuit:battery}{batterij} voedt twee parallelle \omterm{def:g5:series-parallel-circuits:parallel}{takken}: \omterm{def:g5:series-parallel-circuits:parallel}{tak} één draagt
lamp X alleen; \omterm{def:g5:series-parallel-circuits:parallel}{tak} twee draagt de lampen Y en Z \omterm{def:g4:series-circuits:series}{in serie}.
Vergelijk de felheid van de drie lampen, op volgorde, met redenen uit
beide wetten. Voorspel daarna het lot van elke lamp als Z doorbrandt
--- en als in plaats daarvan X doorbrandt.
\end{exercise}

\section{Probleem: het paneel van de vuurtorenwachter}

\begin{problem}[{Weekendprobleem --- de oude vuurtoren opnieuw
bedraden; één batterijbank, vele taken; de storingsnacht van de
wachter}]\label{pb:g7:circuits-series-parallel:1}
De batterijbank van de vuurtoren moet voeden: de grote lamp, een
misthoorn (een machtige \omterm{ex:g6:circuits-and-safety:motor}{zoemer}), het leeslicht van de wachter en
een traplicht --- en de storingsnacht van de wachter overleven.

\medskip\noindent\textbf{Deel I --- Ontwerp.}
\begin{enumerate}
  \item De grote lamp mag nooit zwakker worden doordat er iets
        anders wordt ingeschakeld. Welke bedrading eist dat voor de
        hele installatie, en waarom zou elke koppeling \omterm{def:g4:series-circuits:series}{in serie} de
        schepen verraden?
  \item Elke taak heeft een eigen \omterm{ex:g3:first-electric-circuit:switch}{schakelaar} nodig, en de wachter
        wil één hoofdafsluiting voor stormen. Plaats alle vijf de
        \omterm{ex:g3:first-electric-circuit:switch}{schakelaars}.
  \item De misthoorn en het traplicht mogen nooit door elkaars
        storingen in de weg worden gezeten. Is dat al gegarandeerd?
        Door welke wet?
  \item Schets of beschrijf het volledige schema: bank,
        hoofdschakelaar, vier \omterm{def:g5:series-parallel-circuits:parallel}{takken}, vier
        takschakelaars.
\end{enumerate}

\medskip\noindent\textbf{Deel II --- De vragen van de wachter.}
\begin{enumerate}[resume]
  \item De wachter maakt zich zorgen: ``Vier \omterm{def:g5:series-parallel-circuits:parallel}{taken} die tegelijk
        drinken --- krijgt elke \omterm{def:g5:series-parallel-circuits:parallel}{tak} dan nog steeds de volle duw van
        de bank?'' Antwoord met de parallelwet.
  \item ``En wordt er stroom verbruikt terwijl hij door de grote lamp
        werkt?'' Zet de wachter recht, met de skilift als dat helpt.
  \item Op een mistige nacht draait alles tegelijk. Wat kost de
        parallelle vrijgevigheid, en waar in de vuurtoren wordt dat
        betaald?
  \item Alleen het leeslicht brandt; meetlustig volgt de wachter de
        mars: beschrijf haar volledige rondreis, en noem elk
        onderdeel dat ze passeert en elk dat ze mijdt.
\end{enumerate}

\medskip\noindent\textbf{Deel III --- De storingsnacht.}
\begin{enumerate}[resume]
  \item Middernacht: het traplampje brandt door. Wat gaat er nog
        meer uit? Welke wet geeft het antwoord?
  \item Eén uur: een onhandige reparatie laat een stuk blote draad de
        twee polen van de bank rechtstreeks verbinden,
        stroomopwaarts van elke \omterm{def:g5:series-parallel-circuits:parallel}{tak}. Noem het verschijnsel, en leg
        uit waarom nu elke taak in de vuurtoren tegelijk uitvalt ---
        waar gaat de mars heen?
  \item Twee uur, storing verholpen: de wachter wil de grote lamp
        zo beschermen dat zo'n sluiproute haar nooit meer het zwijgen
        kan opleggen. Er is een tweede, onafhankelijke
        batterijbank beschikbaar. Stel de bedrading voor.
  \item Ochtendgloren: schrijf het logboek van de wachter --- drie
        zinnen: welke wet de schepen vannacht veilig hield, welke
        misvatting de nacht weerlegde, en welke ene draad het ergste
        uur veroorzaakte.
\end{enumerate}
\end{problem}
